% Options for packages loaded elsewhere
\PassOptionsToPackage{unicode}{hyperref}
\PassOptionsToPackage{hyphens}{url}
\PassOptionsToPackage{dvipsnames,svgnames,x11names}{xcolor}
%
\documentclass[
  10pt,
]{article}
\usepackage{amsmath,amssymb}
\usepackage{iftex}
\ifPDFTeX
  \usepackage[T1]{fontenc}
  \usepackage[utf8]{inputenc}
  \usepackage{textcomp} % provide euro and other symbols
\else % if luatex or xetex
  \usepackage{unicode-math} % this also loads fontspec
  \defaultfontfeatures{Scale=MatchLowercase}
  \defaultfontfeatures[\rmfamily]{Ligatures=TeX,Scale=1}
\fi
\usepackage{lmodern}
\ifPDFTeX\else
  % xetex/luatex font selection
\fi
% Use upquote if available, for straight quotes in verbatim environments
\IfFileExists{upquote.sty}{\usepackage{upquote}}{}
\IfFileExists{microtype.sty}{% use microtype if available
  \usepackage[]{microtype}
  \UseMicrotypeSet[protrusion]{basicmath} % disable protrusion for tt fonts
}{}
\makeatletter
\@ifundefined{KOMAClassName}{% if non-KOMA class
  \IfFileExists{parskip.sty}{%
    \usepackage{parskip}
  }{% else
    \setlength{\parindent}{0pt}
    \setlength{\parskip}{6pt plus 2pt minus 1pt}}
}{% if KOMA class
  \KOMAoptions{parskip=half}}
\makeatother
\usepackage{xcolor}
\usepackage[margin=1in]{geometry}
\usepackage{color}
\usepackage{fancyvrb}
\newcommand{\VerbBar}{|}
\newcommand{\VERB}{\Verb[commandchars=\\\{\}]}
\DefineVerbatimEnvironment{Highlighting}{Verbatim}{commandchars=\\\{\}}
% Add ',fontsize=\small' for more characters per line
\newenvironment{Shaded}{}{}
\newcommand{\AlertTok}[1]{\textcolor[rgb]{1.00,0.00,0.00}{\textbf{#1}}}
\newcommand{\AnnotationTok}[1]{\textcolor[rgb]{0.38,0.63,0.69}{\textbf{\textit{#1}}}}
\newcommand{\AttributeTok}[1]{\textcolor[rgb]{0.49,0.56,0.16}{#1}}
\newcommand{\BaseNTok}[1]{\textcolor[rgb]{0.25,0.63,0.44}{#1}}
\newcommand{\BuiltInTok}[1]{\textcolor[rgb]{0.00,0.50,0.00}{#1}}
\newcommand{\CharTok}[1]{\textcolor[rgb]{0.25,0.44,0.63}{#1}}
\newcommand{\CommentTok}[1]{\textcolor[rgb]{0.38,0.63,0.69}{\textit{#1}}}
\newcommand{\CommentVarTok}[1]{\textcolor[rgb]{0.38,0.63,0.69}{\textbf{\textit{#1}}}}
\newcommand{\ConstantTok}[1]{\textcolor[rgb]{0.53,0.00,0.00}{#1}}
\newcommand{\ControlFlowTok}[1]{\textcolor[rgb]{0.00,0.44,0.13}{\textbf{#1}}}
\newcommand{\DataTypeTok}[1]{\textcolor[rgb]{0.56,0.13,0.00}{#1}}
\newcommand{\DecValTok}[1]{\textcolor[rgb]{0.25,0.63,0.44}{#1}}
\newcommand{\DocumentationTok}[1]{\textcolor[rgb]{0.73,0.13,0.13}{\textit{#1}}}
\newcommand{\ErrorTok}[1]{\textcolor[rgb]{1.00,0.00,0.00}{\textbf{#1}}}
\newcommand{\ExtensionTok}[1]{#1}
\newcommand{\FloatTok}[1]{\textcolor[rgb]{0.25,0.63,0.44}{#1}}
\newcommand{\FunctionTok}[1]{\textcolor[rgb]{0.02,0.16,0.49}{#1}}
\newcommand{\ImportTok}[1]{\textcolor[rgb]{0.00,0.50,0.00}{\textbf{#1}}}
\newcommand{\InformationTok}[1]{\textcolor[rgb]{0.38,0.63,0.69}{\textbf{\textit{#1}}}}
\newcommand{\KeywordTok}[1]{\textcolor[rgb]{0.00,0.44,0.13}{\textbf{#1}}}
\newcommand{\NormalTok}[1]{#1}
\newcommand{\OperatorTok}[1]{\textcolor[rgb]{0.40,0.40,0.40}{#1}}
\newcommand{\OtherTok}[1]{\textcolor[rgb]{0.00,0.44,0.13}{#1}}
\newcommand{\PreprocessorTok}[1]{\textcolor[rgb]{0.74,0.48,0.00}{#1}}
\newcommand{\RegionMarkerTok}[1]{#1}
\newcommand{\SpecialCharTok}[1]{\textcolor[rgb]{0.25,0.44,0.63}{#1}}
\newcommand{\SpecialStringTok}[1]{\textcolor[rgb]{0.73,0.40,0.53}{#1}}
\newcommand{\StringTok}[1]{\textcolor[rgb]{0.25,0.44,0.63}{#1}}
\newcommand{\VariableTok}[1]{\textcolor[rgb]{0.10,0.09,0.49}{#1}}
\newcommand{\VerbatimStringTok}[1]{\textcolor[rgb]{0.25,0.44,0.63}{#1}}
\newcommand{\WarningTok}[1]{\textcolor[rgb]{0.38,0.63,0.69}{\textbf{\textit{#1}}}}
\setlength{\emergencystretch}{3em} % prevent overfull lines
\providecommand{\tightlist}{%
  \setlength{\itemsep}{0pt}\setlength{\parskip}{0pt}}
\setcounter{secnumdepth}{-\maxdimen} % remove section numbering
\ifLuaTeX
  \usepackage{selnolig}  % disable illegal ligatures
\fi
\usepackage{bookmark}
\IfFileExists{xurl.sty}{\usepackage{xurl}}{} % add URL line breaks if available
\urlstyle{same}
\hypersetup{
  colorlinks=true,
  linkcolor={Maroon},
  filecolor={Maroon},
  citecolor={Blue},
  urlcolor={blue},
  pdfcreator={LaTeX via pandoc}}

\author{}
\date{}

\begin{document}

\section{DDTA research work plan after documentation-layer
closure}\label{ddta-research-work-plan-after-documentation-layer-closure}

\textbf{WORK PLAN - REVISION 7}

\textbf{Status:} active BA1 execution plan after BA1-T2
split-or-collapse pressure testing.

\textbf{Supersedes:} Revision 6 only for forward execution state; R1-R6
remain historical research records.

\subsection{1. Fixed prior state}\label{fixed-prior-state}

\begin{itemize}
\tightlist
\item
  Chapters 2-4: CLOSED / FINAL for current thesis scope.
\item
  Documentation layer: CLOSED.
\item
  BA0-R systems-modeling prior-art research: CLOSED.
\item
  BA0 responsibility/non-goals: CLOSED by BA0-T3.
\item
  ThreatForge is a case study/tooling instrument, not DDTA semantic
  authority.
\item
  No historical focus set, systems-modeling metaclass or threat-method
  element is inherited automatically into BA1.
\end{itemize}

\subsection{2. BA0 boundary carried
forward}\label{ba0-boundary-carried-forward}

Base Analysis must preserve baseline-scoped shared semantic identity,
source/origin provenance and explicit unresolved state sufficiently for
reusable human/method projections, source drill-down, change-impact
reasoning and source-localized feedback handoff.

The representation remains methodology-neutral and
implementation-independent. Grounded, derived and diagnostic/unresolved
are responsibility-level origin/state distinctions; a generic reviewed
analytical-addition class is not required.

\subsection{3. BA1 objective}\label{ba1-objective}

BA1 decides the \textbf{minimal first-class semantic identity ontology}
needed to realize the BA0 boundary.

A first-class split requires both:

\begin{enumerate}
\def\labelenumi{\arabic{enumi}.}
\tightlist
\item
  independent semantic identity across propositions/projections/change;
  and
\item
  reusable subtype-specific invariants that cannot be represented
  honestly as classifications, roles, values or propositions over an
  existing family.
\end{enumerate}

\begin{Shaded}
\begin{Highlighting}[]
\NormalTok{semantic responsibility exists}
\NormalTok{        !=}
\NormalTok{first{-}class metaclass required}
\end{Highlighting}
\end{Shaded}

\subsection{4. BA1-T1 - minimal ontology candidate
derivation}\label{ba1-t1---minimal-ontology-candidate-derivation}

\textbf{Status: COMPLETED / PROVISIONAL CANDIDATE.}

BA1-T1 produced the initial falsifiable candidate:

\begin{Shaded}
\begin{Highlighting}[]
\NormalTok{BAReferent}
\NormalTok{BAProposition}
\end{Highlighting}
\end{Shaded}

It rejected a fully undifferentiated analytical element and did not
force domain-specific roots.

\subsection{5. BA1-T2 - split-or-collapse pressure
test}\label{ba1-t2---split-or-collapse-pressure-test}

\textbf{Status: COMPLETED / PROVISIONAL PASS WITH MINIMALITY
REFINEMENT.}

BA1-T2 deliberately stressed the two-family candidate against:

\begin{itemize}
\tightlist
\item
  facial-access delivery behavior and specialized constraints;
\item
  order-fulfillment physical handoff reuse;
\item
  lifecycle/correlation-heavy informational meanings;
\item
  internal/external responsibility placement;
\item
  store and contract candidates;
\item
  state/indeterminate semantics;
\item
  a bounded STRIDE-oriented method consumer;
\item
  semantic collapse into one undifferentiated family.
\end{itemize}

\subsubsection{Result}\label{result}

No tested semantic kind forced a dedicated first-class split.

\begin{Shaded}
\begin{Highlighting}[]
\NormalTok{Behavior/Event                    NOT FORCED}
\NormalTok{Information/Resource              NOT FORCED}
\NormalTok{Participant/Component/Capability  NOT FORCED}
\NormalTok{Boundary/Externality              NOT FORCED}
\NormalTok{Store                             NOT FORCED}
\NormalTok{Contract                          NOT FORCED}
\NormalTok{State/Mode/Context                NOT FORCED}
\end{Highlighting}
\end{Shaded}

The semantic one-family collapse was \textbf{REJECTED} because referent
identity and proposition identity/origin/evolution remain independently
required.

\subsubsection{Minimality refinement}\label{minimality-refinement}

R1 allowed \texttt{BAProposition} to be a project-semantic target of
another proposition. T2 found this unnecessary and potentially
misleading.

The refined boundary is:

\begin{itemize}
\tightlist
\item
  project meaning that needs independent reuse/qualification
  -\textgreater{} \texttt{BAReferent};
\item
  methodology-neutral assertion about project meaning -\textgreater{}
  \texttt{BAProposition};
\item
  assertion provenance/diagnostic/change metadata -\textgreater{} BA3
  mechanisms targeting proposition identity.
\end{itemize}

The active candidate is now
\texttt{BA1\_MINIMAL\_BAE\_ONTOLOGY\_CANDIDATE\_R2.md}.

\subsection{6. BA1-T3 - minimal BAE ontology closure
review}\label{ba1-t3---minimal-bae-ontology-closure-review}

\textbf{NEXT / NOT EXECUTED BY THIS PLAN REVISION.}

BA1-T3 is a closure/falsification review, not another ontology invention
exercise.

It must test:

\begin{enumerate}
\def\labelenumi{\arabic{enumi}.}
\tightlist
\item
  \textbf{Necessity:} does removing either \texttt{BAReferent} or
  \texttt{BAProposition} break a closed BA0 responsibility or observed
  corpus requirement?
\item
  \textbf{Sufficiency:} does the two-family boundary cover the tested
  corpora without hiding a still-unresolved first-class identity
  responsibility?
\item
  \textbf{No leakage:} can the boundary be stated without relation
  vocabulary, document types, STRIDE/DFD categories, ThreatForge
  classes, view/tool/storage semantics or general systems-modeling
  taxonomies?
\item
  \textbf{Correct phase placement:} do predicates, roles and actions
  belong to BA2; provenance and identity mechanics to BA3; projections
  to BA4; vocabulary assistance to BA5; and regression to BA6?
\item
  \textbf{Minimality refinement:} is proposition-as-project-target
  correctly excluded from the required core?
\item
  \textbf{Stop criterion:} if no material counterexample remains and
  unresolved details belong later, close BA1 rather than invent BA1-T4
  for reassurance.
\end{enumerate}

\subsubsection{Possible exit}\label{possible-exit}

If the review passes:

\begin{Shaded}
\begin{Highlighting}[]
\NormalTok{BA1 CLOSED}
\NormalTok{BAReferent ACCEPTED first{-}class identity family}
\NormalTok{BAProposition ACCEPTED first{-}class identity family}
\NormalTok{BA2 NEXT}
\end{Highlighting}
\end{Shaded}

If a material counterexample remains, record it explicitly and keep BA1
open. Do not silently push unresolved ontology into BA2.

\subsection{7. Sequence after BA1
closure}\label{sequence-after-ba1-closure}

Do not skip phases:

\begin{Shaded}
\begin{Highlighting}[]
\NormalTok{BA2 {-} Relations and canonical action vocabulary}
\NormalTok{BA3 {-} Document{-}to{-}BAE derivation, provenance and authority}
\NormalTok{BA4 {-} Multi{-}level human and method projections}
\NormalTok{BA5 {-} Controlled vocabulary and optional authoring/extraction assistance}
\NormalTok{BA6 {-} Base Analysis closure and regression}
\end{Highlighting}
\end{Shaded}

\subsubsection{BA2}\label{ba2}

Define the smallest stable relation/action vocabulary over the accepted
BA1 identity ontology. Keep semantic relation, canonical lexical label
and authoring synonym assistance separate.

BA2 may determine proposition arity, participation roles,
dependency/constraint/action semantics and how reusable semantic
classifications are represented, but it must not retroactively create
BA1 root types without a concrete reopen trigger.

\subsubsection{BA3}\label{ba3}

Close identity/equivalence/lifecycle rules, source locators,
grounded/derived/diagnostic materialization, acceptance/rejection, stale
state and document-to-analysis traceability.

\subsubsection{BA4}\label{ba4}

Define derived human and method projections. No view is source
authority.

\subsubsection{BA5}\label{ba5}

Define controlled vocabulary and optional assistance. LLM/NLP proposals
remain candidate-producing support and cannot establish authority
silently.

\subsubsection{BA6}\label{ba6}

Regress the Base Analysis against closed corpora and at least one
structurally different holdout before formal overlay work.

\subsection{8. Analysis envelope remains
downstream}\label{analysis-envelope-remains-downstream}

Only after Base Analysis closure:

\begin{itemize}
\tightlist
\item
  A1 - AnalysisRecord / execution envelope;
\item
  A2 - common reviewed Finding boundary;
\item
  A3 - change/provenance integration;
\item
  formal methodology overlays and STRIDE/STRIDE-AI evaluations.
\end{itemize}

The BA0-T2 STRIDE probe remains only a pre-overlay consumer probe.

\subsection{9. Next authorized
microstep}\label{next-authorized-microstep}

Execute only:

\begin{quote}
\textbf{BA1-T3 - minimal BAE ontology closure review.}
\end{quote}

Do not start BA2, formal STRIDE overlay design, Common Finding schema or
implementation work until BA1 closes explicitly.

\end{document}
